The Jefferson (greatest divisors) method, known internationally as
the d'Hondt method, allocates seats by successively awarding the
next seat to the state or party with the highest priority value
$P_i(s) = p_i/s_i$, where $p_i$ is the state's population (or party's
vote total) and $s_i$ its current seat count.
Equivalently, Jefferson assigns each state $\lfloor p_i/d \rfloor$
seats for a common divisor $d$, always rounding down.

Jefferson was the first apportionment method used in the United States,
employed for congressional apportionments from 1790 to 1840 at
Secretary of State Thomas Jefferson's recommendation.
Its floor rounding systematically underrepresents small states
(whose fractional quota parts are lost to rounding) and
overrepresents large states, a bias that led to its eventual
abandonment in U.S.\ practice.

Internationally, d'Hondt is the dominant proportional representation
formula for party-list parliamentary elections, used in Belgium, Spain,
Portugal, Brazil, Argentina, and dozens of other countries.
We prove the formal equivalence between Jefferson and d'Hondt:
the priority formula $p_i/s_i$ (population-to-seats ratio) and
$V_j/s_j$ (votes-to-seats ratio) are mathematically identical
functions applied to different quantities, confirming that U.S.\
apportionment history and contemporary PR system design address
the same mathematical problem.

We prove that Jefferson maximises total seat share for large states
under any fixed population distribution, characterise the magnitude
of large-state bias for 2020 Census data, and document the historical
1832 controversy that ultimately led to Jefferson's abandonment.
The \textsc{bisect-apportion} crate implements Jefferson/d'Hondt
for both U.S.\ apportionment and comparative electoral systems analysis.
